\begin{homeworkProblem}[Seção 2-16]
  Usando a Segunda Fórmula de Green, conclua que se $u,v$ são funções harmônicas tais que
  \begin{align*}
    \frac{\partial u}{\partial\eta}=\frac{\partial v}{\partial\eta},    
  \end{align*}
  então $u-v=$ constante.
  Em particular, se $u$ é uma função harmônica tal que $\frac{\partial u}{\partial\eta}=0$ em $\partial\Omega$, então $u=$ constante em $\Omega$.
  \begin{solution}
    Note que como $u$ e $v$ são funções harmônicas temos que se definimos $w=u-v$, então $\Delta w=0$, logo, note que pela hipótese temos que
    \begin{align*}
      \frac{\partial w}{\partial \eta}=\frac{\partial u}{\partial \eta}-\frac{\partial v}{\partial \eta}=0.
    \end{align*}
    Pelo que temos que $\frac{\partial w}{\partial \eta}=0$, logo usando a primeira fórmula de Green (que se pode mostrar usando o mesmo  argumento do exercício 2-15), que $\Delta w=0$ e que $\frac{\partial w}{\partial \eta}=0$, temos que
    \begin{align*}
      \int_{\Omega}\left| dw \right|^{2}=&\int_{\Omega}\left\langle dw,dw \right\rangle+w\Delta w
      =\int_{\partial\Omega}\left\langle wdw,\eta \right\rangle
      =\int_{\partial\Omega}w\frac{\partial w}{\partial \eta}
      =0.
    \end{align*}
    Logo, como $\left| dw \right|^{2}\geq 0$, temos que $dw=0$. Depois podemos afirmar que $w=$ constante, mas $w=u-v$, o que nos permite concluir que $u-v=$ constante.\\
    Agora, note que se $v=0$, temos que $\frac{\partial u}{\partial \eta}=0$, logo $u-v=u=$ constante, o que nos permite concluir o exercício. 
  \end{solution}
\end{homeworkProblem}
